\section*{第 \ref{ch:g10:pressure} 章——压强：从运动到潜水}

\begin{solution}{exo:g10:pressure:1}
大接触面：$S = \qty{0.15}{m^2}$，$P = 600/0.15 = \qty{4.0}{kPa}$。
小接触面：$S = \qty{0.060}{m^2}$，$P = \qty{10}{kPa}$。在沙地上，
大接触面\omterm{def:g10:pressure:pressure}{压强}更小，下陷更少。
\end{solution}

\begin{solution}{exo:g10:pressure:2}
$\qty{0.25}{MPa} = \qty{250}{kPa} = \qty{2.5e5}{Pa}$.
$3.0/\num{2.0e-4} = \qty{15}{kPa}$ 超过
$60/0.050 = \qty{1.2}{kPa}$：较小的力\omterm{def:g10:pressure:pressure}{压强}更大。
\end{solution}

\begin{solution}{exo:g10:pressure:3}
$\rho g h = 1000 \times 9.81 \times 3.0 \approx \qty{29}{kPa}$;
\omterm{met:g10:pressure:dive}{绝对压强} $P = 101 + 29 = \qty{130}{kPa}$。
\end{solution}

\begin{solution}{exo:g10:pressure:4}
$V_2 = 100 \times 4.0 / 250 = \qty{1.6}{L}$;
在 \qty{50}{kPa} 下，$V_2 = 400/50 = \qty{8.0}{L}$。
\end{solution}

\begin{solution}{exo:g10:pressure:5}
$F = 65 \times 9.81 \approx \qty{638}{N}$。刀刃：
$638/\num{1.2e-3} \approx \qty{5.3e5}{Pa} = \qty{530}{kPa}$。鞋底：
$638/0.036 \approx \qty{18}{kPa}$。比值 $\approx 30$：刀刃
切入冰中，鞋底则不会。
\end{solution}

\begin{solution}{exo:g10:pressure:6}
拇指：$20/\num{1.2e-4} \approx \qty{1.7e5}{Pa} = \qty{170}{kPa}$。
墙面：$20/\num{1.0e-8} = \qty{2.0e9}{Pa}$。同样的 \qty{20}{N} 力作用于
尖端处的面积小 $12\,000$ 倍，因此只有在那里
\omterm{def:g10:pressure:pressure}{压强}才超过材料所能承受的限度。
\end{solution}

\begin{solution}{exo:g10:pressure:7}
$\Delta P = \rho g h = 1025 \times 9.81 \times 30 \approx
\qty{3.0e5}{Pa}$，因此
$F = \num{3.0e5} \times 0.50 \approx \qty{1.5e5}{N}$——相当于
质量为 $F/g \approx \qty{1.5e4}{kg}$（十五吨）的物体所受的\omterm{def:g10:universal-gravitation:weight}{重力}。
没有船员能逆着海水把舱门向外推开。
\end{solution}

\begin{solution}{exo:g10:pressure:8}
$\rho = \dfrac{P}{g h} = \dfrac{\num{66.2e3}}{9.81 \times 5.00}
\approx \qty{1350}{kg/m^3}$。不是淡水（$\qty{1000}{kg/m^3}$）：
密度明显更大的液体，例如浓盐水。
\end{solution}

\begin{solution}{exo:g10:pressure:9}
$h = P_0/(\rho g)$：水
$\num{1.01e5}/9810 \approx \qty{10.3}{m}$；水银
$\num{1.01e5}/(13600 \times 9.81) \approx \qty{0.76}{m}$。一根
\qty{76}{cm} 的管子放在书桌上即可，而 \qty{10}{m} 的水气压计
需要一个楼梯间——水银的密度占了上风。
\end{solution}

\begin{solution}{exo:g10:pressure:10}
$V = 6000/240 = \qty{25}{mL}$; $P = 6000/75 = \qty{80}{kPa}$.
$P$ 与 $V$ 的关系是双曲线；$P$ 与 $1/V$ 的关系是
过原点的直线，斜率为 \qty{6000}{kPa.mL}。
\end{solution}

\begin{solution}{exo:g10:pressure:11}
在 \qty{20}{m} 深处：$P = 101 + 2 \times 101 = \qty{303}{kPa}$，三
个大气压。在水面处，
$V = 0.50 \times 303/101 = \qty{1.5}{cm^3}$：气泡体积变为原来的三倍。
\end{solution}

\begin{solution}{exo:g10:pressure:12}
\emph{1.} $F = \num{1.01e5} \times \num{1.2e-3} \approx \qty{121}{N}$,
可托住质量 $m = 121/9.81 \approx \qty{12}{kg}$ 的物体。

\emph{2.} $\Delta P = 101 - 20 = \qty{81}{kPa}$:
$F = \num{8.1e4} \times \num{1.2e-3} \approx \qty{97}{N}$，
约能托住 \qty{9.9}{kg} 的物体。
\end{solution}

\begin{solution}{exo:g10:pressure:13}
玻意耳：$12 \times (200 - 50) = \qty{1800}{L}$ 地面\omterm{def:g10:pressure:pressure}{压强}下的空气可用。
在 \qty{30}{m} 深处，每次吸气的\omterm{def:g10:pressure:pressure}{压强}为 \qty{4}{bar}，因此
环境\omterm{def:g10:pressure:pressure}{压强}下 \qty{15}{L/min} 的用气量，
折合地面空气 $15 \times 4 = \qty{60}{L/min}$：$1800/60 =
\qty{30}{min}$。在 \qty{10}{m} 深处（$\qty{2}{bar}$）：
\qty{30}{L/min}，故可用 \qty{60}{min}。深度要用空气来付出代价。
\end{solution}

\begin{solution}{exo:g10:pressure:14}
$P = P_0 \dfrac{V_0}{V} = 101 \times \dfrac{6.0}{1.5} =
\qty{404}{kPa}$。取 $\rho g \approx \qty{10.1}{kPa/m}$，
$h = (404 - 101)/10.1 \approx \qty{30}{m}$.
\end{solution}

\begin{solution}{exo:g10:pressure:15}
\emph{1.} \omterm{def:g10:universal-gravitation:weight}{重力} $F = 1200 \times 9.81 \approx \qty{1.18e4}{N}$；
\omterm{def:g10:pressure:pressure}{压强} $P = F/S = \num{1.18e4}/0.040 \approx \qty{2.9e5}{Pa}$；小
活塞：$f = P s = \num{2.9e5} \times \num{2.0e-4} \approx \qty{59}{N}$
——力被除以 $S/s = 200$。

\emph{2.} 液体体积守恒：
$400 \times 2.0 = \qty{800}{cm^3}$ 必须来自小圆筒，
其行程为 $800/2.0 = \qty{400}{cm} = \qty{4.0}{m}$。功：
一侧为 $59 \times 4.0 \approx \qty{235}{J}$，
另一侧为 $\num{1.18e4} \times 0.020 \approx \qty{235}{J}$——
升降机用距离换取力，而功不变。
\end{solution}

\begin{solution}{pb:g10:pressure:1}
\textbf{1.} $\rho g = 1025 \times 9.81 \approx \qty{10.1}{kPa}$，每
米深度；\qty{10}{m} 海水增加 $\approx \qty{101}{kPa}$，
即每十米约一个大气压，几乎恰好。

\textbf{2.} $P = 101 + 10.1\,h$：\qty{10}{m} 处为 \qty{202}{kPa}，
\qty{20}{m} 处为 \qty{302}{kPa}，\qty{30}{m} 处为 \qty{403}{kPa}，
\qty{40}{m} 处为 \qty{503}{kPa}。

\textbf{3.} $2.0$、$3.0$、$4.0$ 和 $5.0$ 倍 $P_0$。

\textbf{4.} $\Delta P = 403 - 101 = \qty{302}{kPa}$，作用于
$\qty{0.015}{m^2}$ 的面积上：$F = \num{3.02e5} \times 0.015 \approx
\qty{4.5e3}{N}$——相当于面上压着近半吨物体的\omterm{def:g10:universal-gravitation:weight}{重力}。

\textbf{5.} $\Delta P = 10.1 \times 3.0 \approx \qty{30}{kPa}$；
$F = \num{3.0e4} \times \num{6.0e-5} \approx \qty{1.8}{N}$——相当于
一根手指按在鼓膜上的压力，已在泳池深度出现。平衡耳压时
会把空气压入中耳，直到中耳内\omterm{def:g10:pressure:pressure}{压强}与水的\omterm{def:g10:pressure:pressure}{压强}相等。

\textbf{6.} $V = 6.0 \, P_0/P$：$6.0 \times 101/202 = \qty{3.0}{L}$；
$606/302 = \qty{2.0}{L}$；$606/403 = \qty{1.5}{L}$。

\textbf{7.} 当 $P = 2P_0$ 时减半，即在 \qty{10}{m} 处——正处于
下潜的最初一段。

\textbf{8.} 否：$V = 606/P$ 是关于\omterm{met:g10:pressure:dive}{绝对压强}的双曲线
（玻意耳定律关注的是 $P$ 而非 $h$）；每增加 \qty{10}{m}
深度，所减少的体积比前一段更小。

\textbf{9.} 要使 $V = \qty{1.5}{L}$，需满足
$P = 101 \times 6.0/1.5 = \qty{404}{kPa}$，即
$h = 303/10.1 \approx \qty{30}{m}$（见\cref{exo:g10:pressure:14}）。

\textbf{10.} 血液中的血浆转移到胸部的血管中，
血浆不可压缩，占据了空气已不再占据的那部分体积。

\textbf{11.} 环境\omterm{def:g10:pressure:pressure}{压强}为 \qty{403}{kPa}。由玻意耳定律，同一份
空气在 \qty{101}{kPa} 下将占 $4$ 倍的体积，故其
密度是水面空气的 $4.0$ 倍。

\textbf{12.} $V = 6.0 \times 403/202 \approx \qty{12}{L}$——体积翻倍，
仅上升 \qty{20}{m} 便已如此。

\textbf{13.} $V = 6.0 \times 403/101 \approx \qty{24}{L}$：
四倍于他的肺活量。

\textbf{14.} \emph{上升时切勿屏住呼吸。}
只要气道畅通，膨胀的空气便经调节器排出，并且
肺内气体的体积绝不会超过 \qty{6.0}{L}。

\textbf{15.} 莉娜在 \qty{30}{m} 处的 \qty{1.5}{L} \emph{正是}她在水面时
\qty{6.0}{L} 被压缩后的体积：上升时它会重新膨胀至恰好
\qty{6.0}{L}，绝不会更多。马可的 \qty{6.0}{L} 则是在
\qty{403}{kPa} 下\emph{取得}的，其中含有 \qty{24}{L} 的水面空气。

\textbf{16.} $V = 1.0 \times 503/P$：\qty{30}{m} 处为 \qty{1.25}{cm^3}，
\qty{20}{m} 处为 \qty{1.67}{cm^3}，\qty{10}{m} 处为 \qty{2.49}{cm^3}，
水面处为 \qty{4.98}{cm^3}。

\textbf{17.} 各阶段倍率：$503/403 = 1.25$，随后为 $1.33$、$1.49$，
以及最后十米的 $202/101 = 2.0$——增长随
着接近水面而加快。

\textbf{18.} $40/10 = \qty{4.0}{min}$。最后一分钟：从
\qty{10}{m} 上升到水面，\omterm{met:g10:pressure:dive}{绝对压强}减半，是
整个上升过程中相对降幅最大的一段。

\textbf{19.} 上升过快，溶解的氮气会在
血液和关节内形成气泡（即减压病）；缓慢上升，
并在 \qty{5}{m} 附近停留，使气体保持溶解足够长时间，
从而平稳地经肺部排出。

\textbf{20.} 自由潜水者：她的空气只会恢复到
原来的体积——玻意耳定律保护了她。水肺潜水者：在深处取得的空气
上升时成倍膨胀——要呼吸，切勿屏息。而最后十米
会使每一份受困空气的体积翻倍（$\times 2.0$，相比之下，
在 \qty{40}{m} 处下降时仅为 $\times 1.25$）：越接近水面，
上浮就应越慢。
\end{solution}
